\documentclass[10pt,a4paper]{amsart}
\usepackage[margin=19mm]{geometry}
\usepackage{amsmath,amssymb,mathtools}
\usepackage[scr=rsfs,cal=euler]{mathalfa}
\usepackage{xfrac}
\usepackage[unicode,hidelinks,bookmarks=false]{hyperref}
\newcommand{\ie}{i.e.\ }
\newcommand{\cf}{cf.\ }
\newcommand{\resp}{respectively\ }
\newcommand{\Subsection}{\subsection}
\providecommand{\nobreakdash}{\nobreak\hbox{-}}
\setlength{\emergencystretch}{2em}
\multlinegap0pt
\usepackage{amssymb}
\usepackage{algorithm}\usepackage{algpseudocode}
\newtheorem{theorem}{Theorem}
\newcommand{\minimize}{\text{minimize}}
\newcommand{\real}{\mathbb{R}}
\title{Kernel-Based LMI Approaches to \\ Solving the Hamilton–Jacobi–Bellman Equation and \\ Nonlinear Optimal Control}
\author{Boumediene Hamzi, Umesh Vaidya}
\date{Source: arXiv:2603.01084v2 (CC BY 4.0)}
\begin{document}
\maketitle

\noindent\emph{Source and licence.} Kernel-Based LMI Approaches to \\ Solving the Hamilton–Jacobi–Bellman Equation and \\ Nonlinear Optimal Control. Version arXiv:2603.01084v2 (CC BY 4.0), distributed by arXiv under the Creative Commons Attribution 4.0 International licence, http://creativecommons.org/licenses/by/4.0/, as recorded in the arXiv licence metadata for this exact version and shown on the abstract page https://arxiv.org/abs/2603.01084v2. Full licence notice: \texttt{licenses/arxiv-2603.01084-CC-BY-4.0.txt}.\par\smallskip

\noindent\textit{Editorial scope.} Counted unit: the article's theorem thm:collocation (source0.tex lines 363--391). The article's complete original proof of it is delivered with it (source0.tex lines 393--419, one \texttt{proof} environment of 27 lines). No mathematics is written, repaired or re-derived: the statement and the proof are the article's own lines, in the article's own notation, and no retained line is substituted or re-flowed.  No further source-local context is needed: every cross-reference of every retained line resolves inside the two delivered spans.  Nothing was omitted from any retained span, so the package carries no omission record.  No retained line contains a citation, so the wrapper carries no bibliography and no bibliography entry.  No retained line contains a figure environment, an image-inclusion command or drawing code, so no image is delivered and the package declares no asset dependency.  Editorial formatting only, in addition: an AMS \texttt{amsart} preamble that restates the article's own declarations and macros used by the retained lines, namely the environments \texttt{theorem} on separate counters, together with the standard packages those lines need.  The retained environments print in this document as Theorem 1 in order of appearance, whereas the article prints the same items with the numbering of its own full document.  The complete native source of the article as distributed in the arXiv e-print for this exact version is bundled unchanged as \texttt{sources/195537/source0.tex}.
\begin{theorem}[Collocation LMI System with Equilibrium Constraints]
\label{thm:collocation}
Let $\{x_j\}_{j=1}^N \subset \Omega$ be collocation points and $\{x_i\}_{i=1}^M$ be kernel centers. Define the notation:
\begin{align*}
\nabla_k f(x_j) &:= [\nabla_x \kappa(x_j, x_1)^\top f(x_j), \ldots, \nabla_x \kappa(x_j, x_M)^\top f(x_j)]^\top \in \real^M, \\
\nabla_k g(x_j) &:= [\nabla_x \kappa(x_j, x_1)^\top g(x_j), \ldots, \nabla_x \kappa(x_j, x_M)^\top g(x_j)]^\top \in \real^{M \times m}.
\end{align*}

The constrained semidefinite program
\begin{subequations}
\begin{align}
\minimize_{p \in \real^M} \quad & \|p\|^2 \label{eq:sdp_obj} \\
\text{subject to} \quad & \sum_{i=1}^M p_i \kappa(0, x_i) = 0, \label{eq:sdp_boundary} \\
& \sum_{i=1}^M p_i \nabla_x \kappa(0, x_i) = 0, \label{eq:sdp_gradient} \\
& \sum_{i=1}^M p_i \nabla_x^2 \kappa(0, x_i) = P, \label{eq:sdp_hessian} \\
& M_j(p) \succeq 0, \quad j = 1, \ldots, N, \label{eq:sdp_lmi}
\end{align}
\label{eq:sdp}
\end{subequations}
where
\begin{equation}
M_j(p) := \begin{bmatrix}
2(p^\top \nabla_k f(x_j) + q(x_j)) & p^\top \nabla_k g(x_j) \\
(\nabla_k g(x_j))^\top p & D
\end{bmatrix},
\label{eq:Mj_matrix}
\end{equation}
is a convex semidefinite program that prevents the trivial solution while enforcing consistency with linearized optimal control theory.
\end{theorem}

\begin{proof}
We verify that all constraints are convex in $p$:

\textbf{1. Objective function:} $\|p\|^2 = p^\top p$ is convex (quadratic with positive definite Hessian $2I$).

\textbf{2. Boundary constraint \eqref{eq:sdp_boundary}:} This is a linear equality constraint in $p$:
\[
\sum_{i=1}^M p_i \kappa(0, x_i) = k_0^\top p = 0,
\]
where $k_0 = [\kappa(0, x_1), \ldots, \kappa(0, x_M)]^\top$.

\textbf{3. Gradient constraint \eqref{eq:sdp_gradient}:} This consists of $n$ linear equality constraints:
\[
\sum_{i=1}^M p_i \frac{\partial \kappa}{\partial x_\ell}(0, x_i) = 0, \quad \ell = 1, \ldots, n.
\]

\textbf{4. Hessian constraint \eqref{eq:sdp_hessian}:} This consists of $n(n+1)/2$ linear equality constraints (using symmetry):
\[
\sum_{i=1}^M p_i \frac{\partial^2 \kappa}{\partial x_\ell \partial x_k}(0, x_i) = P_{\ell k}, \quad 1 \leq \ell \leq k \leq n.
\]

\textbf{5. LMI constraints \eqref{eq:sdp_lmi}:} The matrix $M_j(p)$ is affine in $p$ because the (1,1) entry $2(p^\top \nabla_k f(x_j) + q(x_j))$ is linear in $p$ plus a constant $2q(x_j)$, the (1,2) entry $p^\top \nabla_k g(x_j)$ is linear in $p$, the (2,1) entry is its transpose and hence also linear, and the (2,2) entry $D$ is constant. An LMI constraint with an affine matrix is convex.

\textbf{6. Preventing trivial solution:} The Riccati constraint \eqref{eq:sdp_hessian} requires $\sum_{i=1}^M p_i \nabla_x^2 \kappa(0, x_i) = P$ with $P \succ 0$. This forces $p \neq 0$, since $p = 0$ would give $0 = P \succ 0$, a contradiction.

Therefore, the optimization problem \eqref{eq:sdp} is a convex semidefinite program.
\end{proof}

\end{document}
